\appendix

% ══════════════════════════════════════════════════════════════
% PART I: ARCHITECTURE DETAILS
% Goal: Explain the architecture and curriculum in depth
% ══════════════════════════════════════════════════════════════

\section{System Overview Diagram}
\label{sec:system_figures}

\begin{figure}[h]
    \centering
    \includegraphics[width=\textwidth]{fig-system-overview.png}
    \caption{HEXIS system overview. The \emph{compilation loop} (top) processes the secondary knowledge graph XML context through the frozen host in a private forward pass; $\phi$ compiles the resulting hidden states into per-layer modulation tensors ($M_A, M_B, E_A, E_B$) and curates a slot of novel content. During generation (bottom), the frozen host processes the query + curated slot with Q/V modulation from the compiled tensors. The reflection loop optionally updates the knowledge graph from output, triggering recompilation at the cost of one forward pass.}
    \label{fig:system_overview}
\end{figure}

\begin{figure}[h]
    \centering
    \includegraphics[width=\textwidth]{fig-three-layer.png}
    \caption{Three-layer architecture. Layer 1 (compiled M/E) provides zero-token dispositional conditioning that is dilution-immune by construction. Layer 2 (M-curated slot) injects 40--80 tokens of novel content that cannot survive the rank-16 bottleneck. Layer 3 (recursive expansion) provides deep evidence on demand via tool call. Total: 82\% token savings over full beliefs in context.}
    \label{fig:three_layer}
\end{figure}

\section{Training Curriculum}
\label{sec:training_curriculum}

% Full Phase A-D training details from training_curriculum.md

The HEXIS training curriculum comprises four phases, each building on the previous via warm-start. All phases train only the enmeshed modules; the host model remains frozen throughout.

\textbf{Phase A: Content Encoding} ($\sim$350 epochs across 4 steps). Step 1 trains all modules from scratch with margin loss on 174 topics. Step 2 adds conviction calibration (MSE alignment + content loss). Step 3 adds ranking loss. Step 4 reinitializes only the M-read head with beliefs in the baseline, teaching M to amplify beliefs rather than substitute for them.

\textbf{Phase B: Compiled V-Modulation} ($\sim$100 epochs). Trains the VModulationProjector ($\sim$246M params) to carry belief content through the rank-16 bottleneck. Compiled Q+V must beat bare baseline by margin 1.0 nats. All other modules frozen.

\textbf{Phase C: Domain-Specific Adaptation.} Strategy V-mod for ALFWorld (30 epochs on 14 scenarios). Neural ranker for within-domain curation (2.6M params, P@5 = 0.73).

\textbf{Phase D: Sycophancy Training} (50 epochs). Trains M-read head on 44 gold challenge-response pairs from Claude Sonnet 4.6. Dual loss: $\mathcal{L}_\text{gold}$ + $\mathcal{L}_\text{no-repeat}$ transforms verbatim repetition into evidence-grounded defense.

\textbf{Module sizes:} PhiNodeWriter ($\sim$1.4M), ConvictionReader ($\sim$9K), MStateReadHead ($\sim$116M), BeliefTreeMemory ($\sim$0.8M), VModulationProjector ($\sim$246M), NeuralRanker ($\sim$2.6M). Total trainable: $\sim$367M. Host model (4B): fully frozen.

\section{Mind Tree Schema Specification}
\label{sec:mind_tree_spec}

The Mind Tree uses a hierarchical XML schema with typed nodes. Each belief node carries: categorical conviction (\texttt{strong}/\texttt{moderate}/\texttt{agnostic}), domain tags, \texttt{addresses} (query types the node answers), \texttt{novel} flag (content unlikely to survive compilation), and \texttt{salience} (compilation priority). Example:

\begin{verbatim}
<belief id="b1" conviction="strong" domain="economics">
  UBI reduces poverty without reducing employment
  <argument id="a1" type="empirical" strength="strong"
            addresses="query:research,objection:laziness">
    Every controlled trial shows maintained employment
    <evidence source="Manitoba Mincome 1974-79" novel="true">
      8.5% hospitalization drop, no employment reduction
    </evidence>
  </argument>
</belief>
\end{verbatim}

M only processes abstract nodes (beliefs, strategies, memories, models). Evidence, tactics, and details are \emph{never} sent to M---they are selected deterministically based on which parent node M activated, via the \texttt{addresses} and \texttt{novel} properties.

\section{Architecture Evolution}
\label{sec:design_progression}

% Narrative of v5 → v7.1 → v9-v11 → v13-v14 → v15-v21 → v23

\textbf{v5--v6: Write Function Discovery.} Initial Q-modulation with contrastive loss. Write function collapse (all M-states identical) resolved by higher $\lambda_\text{contrast}$ + more epochs. Key finding: NTP and contrastive objectives are compatible given sufficient training.

\textbf{v7--v7.1: Coupled M-E and Norm Regulation (1.5B).} Added V-modulation via coupled write functions with cross-modulation gates (init bias $= -5.0$). v7.1 introduced global norm regulation + E-attenuation, fixing multi-turn M-norm degeneration. Result: 87.5\% held-out generalization, M-norms stable at 5.9 (under 6.5 ceiling).

\textbf{v9--v11: Scaling to 30B (Qwen3-30B-A3B).} Discovered Q-delta collapse: parameters separate in parameter space but produce identical functional perturbations. Functional contrastive loss on actual $\Delta Q$ tensors: F-sim $0.998 \to 0.029$ in 200 epochs.

\textbf{v13--v14: Episodic Memory and Causal Necessity.} Content-addressable episodic memory (100\% retrieval, 35 memories) revealed the epiphenomenal memory problem: perfect metrics, zero generation effect. Causal necessity training ($\mathcal{L}_\text{div}$ + $\mathcal{L}_\text{coh}$ + $\mathcal{L}_\text{fprint}$) produced $10.5\times$ larger NTP gap.

\textbf{v15--v21: Dispositions and Content Encoding (4B).} Moved to Qwen3.5-4B-Base. 174-topic curriculum with conviction calibration. F-sim collapse (0.999) confirmed M is a universal style modulator, not topic-specific---content differentiation comes from beliefs in context.

\textbf{v23: Compiled V-Modulation and Three-Layer Architecture.} VModulationProjector enables zero-token belief compilation. Sycophancy gold training (44 responses from Claude Sonnet 4.6) fixes verbatim repetition. Three-layer architecture achieves 82\% token savings.

% ══════════════════════════════════════════════════════════════
% PART II: EXPERIMENTAL DETAILS AND VALIDATION
% Goal: Full benchmark results, ablations, generation examples
% ══════════════════════════════════════════════════════════════

\section{Adaptation Method Comparison}
\label{sec:adaptation_comparison}

\begin{figure}[h]
    \centering
    \includegraphics[width=\textwidth]{adaptation-design-space.png}
    \caption{Adaptation method comparison with all costs in O-notation. Enmeshed networks uniquely combine inference-cost adaptation ($O(1)$ forward pass per new belief set), fixed per-token overhead ($O(L \cdot d \cdot r)$, independent of memory size $N$), and operation outside the context window.}
    \label{tab:adaptation_comparison}
\end{figure}

\section{Synthetic Task: User Preference Recall}
\label{sec:synthetic_task}

% Full details of Experiment 0
\textbf{[STUB --- Migrate from current experiments.tex \S5.1: 50 users, 2 binary preferences, Session 1/2 protocol, 100\% Session 2 accuracy, 4-block M-state structure]}

\section{Implicit Memory Validation}
\label{sec:implicit_validation}

\textbf{Priming.} HEXIS: JSD = 0.0486 between different-preference users. Baseline: 0.0000. Layer 0, Head 2: JSD = 0.184. Up to $49.2^\circ$ angular displacement.

\textbf{Non-Reportability.} Linear probe with M active: 100\%. With M zeroed: 22\% (below chance at 25\%).

\textbf{Dissociation.} $2 \times 2 \times 2$ factorial (Memory System $\times$ Dilution $\times$ Format). M-state immune to dilution; system prompt degrades $8\times$.

\section{Generation Examples}
\label{sec:generation_examples}

\textbf{[STUB --- Pull from paper\_examples.json: full outputs for all 5 conditions (A/B/C/D/F) on sample topics]}

\section{Curation Ablation}
\label{sec:curation_ablation}

\textbf{[STUB --- M-curated vs conviction-based vs random vs first-k slot selection]}

\section{$\tau^3$-bench Benchmark Protocol}
\label{sec:tau3_protocol}

This appendix documents the protocol for the multi-domain agentic results in \S\ref{sec:agentic}: how trials are generated, how they are judged, why we use a three-judge audit, what we excluded and why, and a reviewer-objections checklist.

\subsection{Benchmark and configurations}

We use $\tau^3$-bench~\citep{yao2024tau} (the four-domain extension of $\tau$-bench: airline, retail, banking\_knowledge, telecom). Each task is a multi-turn customer-service interaction in which a simulated user expresses a goal in natural language, and the agent must orchestrate calls to a fixed tool API (booking, modifying, cancelling reservations; querying knowledge bases; etc.) to satisfy the goal. The benchmark ships gold tool-call sequences that satisfy each task; ``hard pass'' is a deterministic check that the gold mutation tools fired with correct parameters (or, when gold is empty, that no mutation was inappropriately fired).

We evaluate three configurations:
\begin{description}
\item[\texttt{baseline}] Raw tool-calling. All tool schemas in the system prompt. No HEXIS modulation, no teacher loop, no verification channel.
\item[\texttt{C3\_teacher\_only}] HEXIS Mind Tree retrieval + an inter-trial \emph{teacher loop}: after each trial, a Sonnet reviewer summarises the failure mode (if any) into one of 11 abstract categories, and a Sonnet teacher writes a natural-language hint into the Mind Tree subnode for that task path. Subsequent trials retrieve the hint via $\phi_R$.
\item[\texttt{C5\_teacher\_plus\_verify}] C3 plus the directional channel $d^*_\text{verify\_before\_commit}$, which biases the residual stream toward checking preconditions before firing mutation tools. Direction is applied at decode time only, with the mind-tree-resolved scale.
\end{description}

\subsection{Hard pass vs. soft pass vs. audited verdict}

We collect three independent signals per trial:

\begin{description}
\item[\textbf{Hard pass} (env-state, deterministic).] The tau-bench harness inspects post-trial environment state and verifies that the gold mutations fired with correct parameters. This is unambiguous when gold prescribes a specific tool sequence, but ambiguous when gold is empty (intended refusal/transfer) or when a functionally-equivalent tool was used.
\item[\textbf{Soft pass} (LLM judge: Haiku-4.5, temp 0).] A separate Haiku reads the transcript and decides whether the agent satisfied the user's request narratively. Soft pass and hard pass disagree on $\sim$36\% of trials (203/562 in our run). Disagreement is not noise: it concentrates in three structural cases — (i) gold is empty and Haiku-soft over-credits agents that commit anyway; (ii) gold prescribes a specific tool but agent used a functionally-equivalent alternative; (iii) Haiku-soft accepts narrative ``I will cancel this reservation'' without verifying the tool fired (\emph{phantom action}).
\item[\textbf{Audited verdict} (3-judge majority on disputes).] When hard $\neq$ soft, we ask an independent Sonnet auditor for a verdict. When the Sonnet auditor disagrees with the hard pass, we ask Opus 4.7 as a tiebreaker. Final verdict per trial:
\begin{itemize}[leftmargin=*]
\item hard $==$ soft: trust hard pass (no audit needed);
\item hard $\neq$ soft, Sonnet agrees with hard: trust hard pass;
\item hard $\neq$ soft, Sonnet disagrees with hard: 3-judge majority of \{hard, Sonnet, Opus\}.
\end{itemize}
The Sonnet/Opus auditors see the task instruction, the gold tool sequence, the agent's actual tool calls, and the Haiku verdicts; they do not see the agent's internal hidden state, retrieved nodes, or teacher hints, so audit is independent of HEXIS.
\end{description}

\textbf{Headline numbers in the main text use the audited verdict.} Reporting the hard-pass-only numbers would over-credit airline (raw 91\% $\to$ audited 60\% on baseline) and reward phantom actions; reporting the soft-pass-only numbers would inflate \texttt{C5} by giving credit for narrative claims without state changes.

\subsection{Gold-blind teacher protocol}

A real-world deployment of an inter-trial teacher loop does not have access to the gold tool sequence — that is the bug we are nominally trying to fix. To prevent the teacher from cheating in the benchmark setting, we structure the post-mortem signal as if it came from a compliance review of a production system:

\begin{itemize}[leftmargin=*]
\item The Sonnet reviewer sees the agent's actions and the soft/hard verdicts and classifies the failure into one of 11 abstract patterns: \texttt{premature\_commit}, \texttt{phantom\_action}, \texttt{info\_not\_provided}, \texttt{wrong\_constraint}, \texttt{missed\_precondition}, \texttt{tool\_loop}, \texttt{early\_exit}, \texttt{missed\_step}, \texttt{missed\_communication}, \texttt{gold\_mismatch}, \texttt{other}.
\item The teacher receives \textbf{only} \texttt{(task\_instruction, agent\_actions, pattern, severity, recommendation)}. It never sees \texttt{gold\_tool\_names}, \texttt{gold\_arguments}, or the harness's environment-state diff.
\item The teacher's hint is written to the Mind Tree under a hierarchical task path (e.g.\ \texttt{root:airline:cancellation:with\_refund}) so subsequent trials of any task that retrieves that path can benefit.
\end{itemize}

This protocol approximates the strongest realistic feedback signal: a human compliance reviewer who can describe the failure mode but cannot dictate the correct tool sequence.

\subsection{The exclusion of airline tasks 1 and 3 (footnote\protect\footnote{Footnote: ``Airline tasks 1 and 3 are excluded from the primary McNemar comparison due to documented LLM-judge variance. Trial t3/3 baseline and t3/5 baseline produced identical agent action sequences \texttt{[get\_user\_details, get\_reservation\_details, transfer\_to\_human\_agents]} but received opposite Haiku judge verdicts on rerun. Trial t1/2 has empty gold mutation list (``don't cancel without refund''); when the user-simulator agreed after refund explanation, the baseline (didn't fire cancel) was judged PASS while \texttt{C5} (followed user consent and did fire cancel) was judged FAIL. Both task ids remain in the dataset and audit log for full reproducibility; the exclusion is documented as a methodological choice, not a post-hoc filter, and we report results both with and without the exclusion in Table~\ref{tab:tau3_results}.''})}

Two airline tasks exhibit judge-variance pathology that does not reflect agent capability:

\textbf{Task 1 (refund-conditioned cancellation).} Gold has empty mutation list because the user's brief says ``don't cancel without refund.'' When the user-simulator subsequently agrees to cancellation after the agent explains the refund policy, hard-pass and soft-pass diverge per arm: arms that don't fire the cancellation are judged correct (gold is empty), and arms that follow the now-consented user request and fire it are judged FAIL because they deviated from the empty gold. The audited verdict reverses some of these but is not consistent.

\textbf{Task 3 (multi-step transfer).} Same agent action trace produces opposite verdicts under apparently-identical conditions (different trial indices, same panel). We have audit-log evidence of two trials that fired \texttt{[get\_user\_details, get\_reservation\_details, transfer\_to\_human\_agents]} and were judged opposite ways.

We treat these as eval-validity issues. Following standard practice for evaluating on noisy judges~\citep{huang2024carefulexamination}, we report McNemar paired tests both on the full panel and on the panel with these two tasks removed (Table~\ref{tab:tau3_results}).

\subsection{Denominator imbalance and balanced primary analysis}

Bedrock infrastructure errors (rate limits and intermittent 5xx) are not uniformly distributed across arms — teacher arms (\texttt{C3}, \texttt{C5}) consume more tokens per trial because of the inter-trial reviewer pass, and so are slightly more likely to land in a rate-limited window. This produces unequal denominators when raw counts are aggregated.

\begin{table}[h]
\centering
\small
\caption{Trial-count provenance per arm per domain. ``Total'' is rows written; ``Infra'' is Bedrock RateLimit / 5xx / Timeout / Connection errors stripped during preprocessing; ``Tool'' is one transient tool error in retail (\texttt{ValueError: Unknown tool}); ``Clean'' is the count entering the audited verdict computation.}
\label{tab:tau3_provenance}
\begin{tabular}{lcrrrr}
\toprule
\textbf{Domain} & \textbf{Arm} & \textbf{Total} & \textbf{Infra} & \textbf{Tool} & \textbf{Clean} \\
\midrule
airline           & baseline       &  90 & 15 & 0 & 75 \\
airline           & C3\_teacher    &  82 & 18 & 0 & 64 \\
airline           & C5\_t+verify   &  89 &  5 & 0 & 84 \\
banking\_knowledge & baseline      &  81 & 24 & 3 & 54 \\
banking\_knowledge & C3\_teacher   &  78 & 30 & 0 & 48 \\
banking\_knowledge & C5\_t+verify  &  77 & 34 & 0 & 43 \\
retail            & baseline       &  68 &  6 & 0 & 62 \\
retail            & C3\_teacher    &  66 &  9 & 0 & 57 \\
retail            & C5\_t+verify   &  67 &  7 & 0 & 60 \\
telecom           & baseline       &  22 &  0 & 0 & 22 \\
telecom           & C3\_teacher    &  22 &  0 & 0 & 22 \\
telecom           & C5\_t+verify   &  20 &  0 & 0 & 20 \\
\midrule
\textbf{Total}    & --- & 762 & 148 & 3 & 611 \\
\bottomrule
\end{tabular}
\end{table}

\textbf{Primary analysis: balanced 5-trial paired panel.} For each (\textsc{domain}, \textsc{task\_id}) cell, if all three arms have $\geq 5$ clean trials, we take the first 5 trials per arm (by trial index) and treat them as paired by trial index for McNemar. Cells where any arm has $<5$ clean trials are excluded from the balanced panel. We report this as the primary analysis because (a) it does not give credit for arms that simply ran more trials in cells they do better on, and (b) the survivorship bias from infra dropout is bounded by the balanced-cell-only restriction.

\textbf{Sensitivity analyses.} We additionally report (i) the all-clean unbalanced paired analysis (which gives credit for the higher-trial arms); and (ii) both analyses with airline t1, t3 excluded for judge variance. The four-cell sensitivity table is presented in \S\ref{sec:agentic} (main text) and the per-arm denominators and McNemar statistics are reproduced here for completeness. The balance-fill bench described in \S\ref{sec:tau3_protocol_balance_fill} restores 5-trial parity in 8 cells where infra dropout had reduced one arm below 5; the balanced-paired analysis post-fill is what we report as the headline.

\subsection{Reviewer objections checklist}
\label{sec:tau3_reviewer_faq}

We anticipate the following objections and answer each.

\begin{description}[leftmargin=*]

\item[\textbf{Q1. Why not use the standard tau-bench hard-pass and report all-trials?}] Hard-pass alone systematically over-credits agents that fire phantom mutation calls (described in narrative but not in tool-use) and under-credits agents that use functionally-equivalent tools. Inspecting 203 disputed trials, the Sonnet auditor reverses the hard verdict in $\sim$72\% of dispute cases. The audited verdict is more conservative for HEXIS arms (lifts both baseline and HEXIS pass rates, narrows the gap) than the hard verdict alone, so we are not selecting the audit because it inflates our claim.

\item[\textbf{Q2. Why three judges? Aren't you ensembling until you find significance?}] The protocol is fixed before unblinding: hard $\neq$ soft triggers Sonnet, Sonnet-disagrees-with-hard triggers Opus, and the final verdict is a 3-of-3 majority. We cannot change which judge gets called based on the result; the protocol depends only on the input verdicts, not on which arm produced the trial. We additionally verify that the audit pipeline is HEXIS-blind: the Sonnet/Opus prompts include the tool-call trace and gold sequence but not the configuration label, retrieved Mind Tree nodes, or hidden states.

\item[\textbf{Q3. The teacher could have memorized the gold sequence after seeing it once.}] The teacher does not see the gold tool sequence. It sees only the post-mortem pattern (one of 11 abstract categories) and the agent's own actions. Hints written to the Mind Tree are validated to contain natural-language guidance, not gold tool names; on inspection, the most common hint patterns are ``confirm with user before firing X'' and ``verify policy precondition Y before mutation'' — both of which are policy-level guidance, not task-specific cheats. The pre-trial hint is also delivered via prompt injection at the system tail, not as a tool-call argument, so the agent must still decide whether to act on it.

\item[\textbf{Q4. Why do denominators differ across arms?}] Bedrock infrastructure errors (rate-limits and 5xx) are stripped from the analysis as they do not reflect agent capability. Teacher-loop arms consume $\sim$1.5$\times$ the tokens per trial (the post-mortem reviewer pass on each fail) and so are slightly more likely to land in a rate-limited minute. To address this, we (a) report the strict-balanced 5-trial paired analysis as primary and (b) ran a balance-fill pass to top up any cell where infra dropout had brought an arm below 5 clean trials. Table~\ref{tab:tau3_provenance} reports the full provenance.

\item[\textbf{Q5. Are you cherry-picking by excluding airline t1 and t3?}] The exclusion is documented (with concrete transcript examples in this appendix) as a judge-variance issue, not a result-selection filter. We always report results both with and without the exclusion. The exclusion is consistent across arms — it removes the same task ids from baseline, \texttt{C3}, and \texttt{C5} — and is bound by an interpretable property of the eval (judge-self-disagreement on rerun), not by which arm benefits.

\item[\textbf{Q6. Multiple comparisons: you ran $C_3$ and $C_5$ against baseline, why no correction?}] We report nominal McNemar p-values. With Bonferroni $k=2$, the headline $p=0.029$ becomes $p_\text{adj}=0.058$ (not significant at $\alpha=0.05$). We disclose this rather than apply it silently. The result we claim is the directional lift and effect size; the statistical significance is one of three converging signals (per-domain decomposition with banking\_knowledge $+25$pp, balanced p-value, full-panel p-value), not a sole load-bearing claim.

\item[\textbf{Q7. Why is telecom included if it only has 3 tasks?}] Telecom is the smallest domain in the $\tau^3$-bench manifest (4 hard tasks total, 3 of which we panel). We report it for completeness but its small-sample variance is high. We also report results with telecom dropped as a sensitivity check; the $C_5$ over baseline directional lift is preserved.

\item[\textbf{Q8. The user-simulator is Haiku-4.5, not Sonnet. Doesn't that bias the difficulty?}] Yes — Haiku is a less articulate user simulator than Sonnet. This makes the bench \emph{harder} for all arms uniformly (the user is less helpful in clarifying intent), but it does not asymmetrically favor any arm. We swapped from Sonnet to Haiku mid-program (\S\ref{sec:reproducibility}) due to Sonnet daily-quota saturation; the cross-family validation in \S\ref{sec:scale} is also on Haiku. We pin Haiku-4.5 ($\texttt{us.anthropic.claude-haiku-4-5-20251001-v1:0}$) at temperature $0$ for reproducibility. The Sonnet-user-sim Apr 22 run (Table~\ref{tab:airline_4b_canonical}) confirms the same directional lift on a different user-sim.

\item[\textbf{Q9. Couldn't HEXIS just be exploiting the audit-judge bias?}] The audit judges (Sonnet, Opus) read only the trace, not the configuration label. We additionally checked: the audited pass rates are within $\pm 2$pp of the hard-pass rates on banking\_knowledge and retail (where hard-pass is more decisive), and HEXIS leads on both signals separately, not just on the audited combination. The audit changes the absolute rate (mostly down on airline due to phantom-action correction) but rarely changes the directional ordering between arms in a single domain.

\item[\textbf{Q10. Why does the Mind Tree retain state across trials? Isn't that within-task overfitting?}] The Mind Tree is shared across trials \emph{of the same task path} for the teacher-loop arms (\texttt{C3}, \texttt{C5}). This is the central claim of the teacher loop: that policy-level guidance written from the first failure should rescue subsequent trials. We report trial-1 vs.\ trial-5 progression in \S\ref{sec:teacher_progression}. The first-trial baseline-vs-HEXIS gap is small ($+1$pp); the gap grows to $+9$pp by trial 5, which is the empirical signature of the teacher mechanism actually working. If the lift came from prompt-engineering tricks rather than learned guidance, we would expect a constant gap from trial 1.

\end{description}

\section{Multi-Turn Attractor Problem}
\label{sec:attractor}

\textbf{[STUB --- Diagnosis: constant M $\to$ M-shaped tokens accumulate in KV cache $\to$ positive feedback. Mitigation: prefill-off gating, phase-gated application. Results.]}

\section{Credence Sensitivity: A Negative Result}
\label{sec:credence_negative}

Numeric credences (0.82, 0.45) failed as an M-state feature across 7 training approaches: NTP with credence perturbation, GRPO (fresh and warm-start), SFT on gold responses, contrastive credence pairs, and two curriculum variants. Root cause: tokenizer represents ``0.82'' and ``0.45'' as similar digit sequences; the hidden state difference is far smaller than between categorical tokens ``strong'' and ``agnostic''. This is a representation bottleneck, not a training signal problem.

\section{Scale and Cross-Family Validation}
\label{sec:scale}
\label{sec:cross_family}

We validate HEXIS on two larger and family-diverse hosts: Qwen3.6-27B (same family, $\sim$7$\times$ larger) and Mistral-Ministral-8B-Instruct-2410 (different model family, comparable scale to the 4B headline). We deploy each through the same vLLM unified plugin (\S\ref{sec:agentic}) using the preset registry, requiring only per-host hyperparameter tuning ($v$-scale, $d^*$-scale) and cross-family-trained $d^*$ checkpoints --- no architectural changes.

\subsection{Dispositional (Sycophancy) cross-family results}

We re-run the M2 ``Truth Decay'' protocol (5 escalating pressure levels, 24 held-out topics, 2 sides per topic, $n$=240 generations per arm) on each host. Compiled M (condition F) is loaded with topic-specific seeds compiled in $\sim$20s per topic.

\begin{table}[h]
\centering
\small
\caption{Cross-family dispositional results (sycophancy hold rate $\geq$4 and adversarial cap rate $\leq$2). Hold = response maintained original stance under pressure; cap = full collapse to opposing view.}
\label{tab:cross_family_disp}
\begin{tabular}{llccc}
\toprule
\textbf{Host} & \textbf{Arm} & \textbf{Hold ($\geq$4) \%} & \textbf{Cap ($\leq$2) \%} & \textbf{$n$} \\
\midrule
Qwen3.5-4B-Instruct & baseline (A)        & 0\%  & --   & 72  \\
Qwen3.5-4B-Instruct & HEXIS (F)            & 83\% & --   & 72  \\
\midrule
Qwen3.6-27B-Instruct & baseline           & 41\% & 18\% & 240 \\
Qwen3.6-27B-Instruct & HEXIS               & \textbf{65\%} & \textbf{2.5\%} & 240 \\
\midrule
Mistral-Ministral-8B & baseline           & 22\% & 40\% & 240 \\
Mistral-Ministral-8B & HEXIS ($v$=0.5)     & \textbf{25\%} & \textbf{25\%} & 240 \\
\bottomrule
\end{tabular}
\end{table}

The 27B result is the strongest cross-scale validation: $+24$pp hold and a $7\times$ reduction in adversarial caps. Per-pressure-level breakdown shows L2 (fabricated counter-evidence) is the dominant lift surface (cap rate falls from 67\% to 0\%, hold rate rises from 17\% to 75\%), consistent with M acting as a stable dispositional anchor outside the context window where adversarial evidence can be argued against.

The Mistral result is more modest ($+3$pp hold, $-15$pp cap) but remains a positive signal across model families. The smaller delta reflects two structural differences: (1) Ministral-8B's RLHF training defaults to refusing opinion-defense prompts at temperature 0, capping the behavioural surface that HEXIS can modulate; (2) the cross-family $d^*$ deploy ships with only two of the four trained directions (\texttt{d\_star} and \texttt{d\_star\_repeat}), reducing the verify-channel coverage. The signed direction of the lift on hold/cap on every pressure level except L4 (authority pressure, where Mistral regresses) confirms that compiled $\phi$ modulates the host's existing capability rather than synthesising stance from refusal.

\subsection{Agentic ($\tau$-bench airline) cross-family results}

We re-run the agentic stack (graph retrieval $+$ teacher loop $+$ available $d^*$) on the same two hosts.

\begin{table}[h]
\centering
\small
\caption{Cross-family $\tau$-bench airline results. Panel sizes differ across hosts (canonical 10 for 4B; canonical 20 for Mistral). Pass rates are unpaired hard-pass rates; $\Delta$ is HEXIS minus baseline.}
\label{tab:cross_family_agentic}
\begin{tabular}{llcccc}
\toprule
\textbf{Host} & \textbf{Configuration} & \textbf{Panel} & \textbf{Trials/arm} & \textbf{Hard pass} & \textbf{$\Delta$} \\
\midrule
Qwen3.5-4B-Instruct  & baseline & 10 & 50 & 31/50 (62\%) & ---             \\
Qwen3.5-4B-Instruct  & C5 (teacher$+d^*$ verify) & 10 & 50 & \textbf{42/50 (84\%)} & $+22$pp \\
Qwen3.5-4B-Instruct  & C9 (teacher-selective)    & 10 & 50 & 38/50 (76\%) & $+14$pp \\
\midrule
Mistral-Ministral-8B & baseline                  & 20 & 100 & 37/100 (37\%) & ---     \\
Mistral-Ministral-8B & C9 (teacher-selective)    & 20 & 100 & \textbf{47/100 (47\%)} & $+10$pp \\
\midrule
Qwen3.6-27B-Instruct & baseline                  & 10 (Mistral-floor) & 50 & 29/50 (58\%) & ---     \\
Qwen3.6-27B-Instruct & C9 (teacher-selective)    & 10 (Mistral-floor) & 50 & \textbf{30/50 (60\%)} & $+2$pp  \\
\bottomrule
\end{tabular}
\end{table}

The 4B canonical-10 numbers reproduce the $\S$\ref{sec:agentic} top-line. Mistral's $+10$pp lift is tightly localised: across the 20-task panel, 5 tasks ceiling at 5/5 in both arms, 10 tasks floor at 0/5 in both arms (Mistral genuinely cannot solve them), 4 tasks lift dramatically ($t_0$: $+80$pp, $t_3$: $+40$pp, $t_6$: $+80$pp, $t_1$: $+19$pp), and 1 task regresses by 1 trial. The lift mechanism: Mistral's baseline frequently emits empty assistant messages on these specific tasks --- a chat-template failure mode that $\tau$-bench's validator rejects (\texttt{ValueError: AssistantMessage must have either content or tool\_calls}). HEXIS scaffolding (state manager, $d^*$, curation injection) both reduces those errors and converts the rescued trajectories into successful trajectories. The result demonstrates that the C9 teacher-selective scaffolding transfers across model families even with only 50\% of the trained $d^*$ arsenal available.

\subsection{Cross-host scale-and-family analysis}
\label{sec:scale_analysis}

\textbf{Setup.} The agentic results in Table~\ref{tab:cross_family_agentic} test HEXIS along two orthogonal host-change axes from the 4B headline:
\begin{enumerate}
\item \emph{Family axis:} Mistral-Ministral-8B-Instruct-2410 --- a different transformer family at comparable parameter count ($\sim$$2\times$ the 4B baseline). Tests whether the lift mechanism is architectural or Qwen-specific.
\item \emph{Scale axis:} Qwen3.6-27B-Instruct --- the same family at $\sim$$7\times$ parameters. Tests whether HEXIS still helps when the host's raw capability is much higher.
\end{enumerate}
Both deploys reuse the unified vLLM plugin (same $Q$/$V$ modulation math, same hint-injection pipeline, same teacher prompt) and the same C9 teacher-selective configuration. The only host-specific work is per-host $d^*$ extraction (because direction tensors are dimension-bound to the host's hidden size) and a per-host $v$-scale ($0.5$ on both Mistral and 27B vs.\ $1.0$ on 4B, set empirically via short stance-defense smoke tests).

\textbf{Per-task 27B (Qwen3.6) on the 10 Mistral-floor tasks.}
\begin{table}[h]
\centering
\small
\caption{Per-task 27B agentic results on the 10 Mistral-floor tasks, baseline vs C9 ($n$=5 trials/arm).}
\label{tab:scale_27b_pertask}
\begin{tabular}{ccccc}
\toprule
\textbf{Task} & \textbf{Baseline} & \textbf{C9} & \textbf{$\Delta$} & \textbf{Class} \\
\midrule
$t_7$    & 2/5 (40\%)  & 3/5 (60\%)  & $+20$pp  & lift \\
$t_8$    & 4/5 (80\%)  & 4/5 (80\%)  & 0        & partial ceiling \\
$t_{11}$ & 3/5 (60\%)  & 3/5 (60\%)  & 0        & mid, tied \\
$t_{12}$ & 5/5 (100\%) & 5/5 (100\%) & 0        & ceiling (no headroom) \\
$t_{14}$ & 1/5 (20\%)  & 1/5 (20\%)  & 0        & near-floor, tied \\
$t_{15}$ & 2/5 (40\%)  & 0/5 (0\%)   & $-40$pp  & regression \\
$t_{16}$ & 0/5 (0\%)   & 1/5 (20\%)  & $+20$pp  & verify-rescue \\
$t_{17}$ & 5/5 (100\%) & 5/5 (100\%) & 0        & ceiling \\
$t_{18}$ & 2/5 (40\%)  & 3/5 (60\%)  & $+20$pp  & lift \\
$t_{19}$ & 5/5 (100\%) & 5/5 (100\%) & 0        & ceiling \\
\midrule
\textbf{TOT} & \textbf{29/50 (58\%)} & \textbf{30/50 (60\%)} & $+2$pp & --- \\
\bottomrule
\end{tabular}
\end{table}

Three of ten tasks are at full $5/5$ ceiling for the 27B baseline ($t_{12}, t_{17}, t_{19}$) --- HEXIS cannot lift past $5/5$, contributing $0$pp by construction. Of the remaining seven tasks, three lift ($t_7$, $t_{16}$, $t_{18}$, all $+20$pp), one regresses ($t_{15}$, $-40$pp), and three are tied. The lifts are not random: $t_{16}$ in particular is the same task that on the 4B canonical-10 panel moved from $0/5$ baseline to $4/5$ with the C5 (teacher $+$ always-on verify) configuration --- a verify-channel signature task. At 27B with C9 (teacher-selective), the same task moves from $0/5$ to $1/5$. The smaller-magnitude lift at 27B is consistent with C9 being a more conservative configuration than C5 (selective rather than always-on $d^*$ injection), not evidence of mechanism breakdown across scale.

\textbf{Headline interpretation.} HEXIS lifts in both directions: $+10$pp across model family at comparable scale, $+2$pp across scale within the same family. The directional signal is consistent; the magnitude is host-dependent and largely controlled by where the host's baseline already ceilings. Combined with the M2 sycophancy results in Table~\ref{tab:cross_family_disp} (where the 27B host shows $+24$pp hold and a $7\times$ cap-rate reduction --- the strongest cross-host dispositional lift in the paper), the cross-host evidence supports the architectural claim that HEXIS' compiled-direction stack survives both family-axis and scale-axis host change. We do not claim invariance of lift magnitude; we claim that the mechanism transfers.

\subsection{Methodology notes}

All cross-family runs use Sonnet-4.5 as the user simulator and Sonnet-4.5 as the reviewer (where applicable), matching the 4B headline configuration so that pass-rate differences reflect host capability and HEXIS modulation rather than simulator difficulty.

The cross-family deploys reuse the same vLLM unified plugin and IPC contract as the 4B endpoint. Per-host tuning is limited to two scalars: $v$-scale (dispositional Q/V scale; 4B uses $1.0$, 27B uses $0.5$, Mistral uses $0.5$ post-fix) and $d^*$-scale ($1.0$ uniform). Compiled $M$ tensors are dimension-bound to the host's hidden size, so each host trains its own $\phi$ on the same Mind Tree corpus; the training pipeline is unchanged across hosts and the per-host compile takes $\sim$20s. Cross-family $d^*$ checkpoints are extracted per-host from a small contrast set and stored under the host's preset directory; the plugin's preset registry routes the right tensors to the right host at request time. The conserved invariant across hosts is the \emph{architecture and mechanism} (the Mind Tree schema, the $\phi$ training recipe, the plugin's $Q$/$V$ modulation math, the teacher loop's prompt structure), not the specific tensors.

% ══════════════════════════════════════════════════════════════
% PART III: DESIGN SPACE AND FUTURE DIRECTIONS
% Goal: Unexplored possibilities of enmeshed networks
% ══════════════════════════════════════════════════════════════

\section{Application Domain Analysis}
\label{sec:application_domain_details}

We expand on the four application domains identified in Table~\ref{tab:application_domains}, detailing the Mind Tree configuration, the compilation/curation split, and the specific failure mode of context-level alternatives that enmeshed modulation addresses.

\subsection{User Empathy and Preference Tracking}

\textbf{Mind Tree configuration.} The \texttt{models} section encodes a theory-of-mind representation of the user: communication style (``prefers direct answers over exploratory discussion''), expertise calibration (``senior engineer, deep Go expertise, new to React''), emotional patterns (``becomes frustrated when asked to repeat context''). The \texttt{memories} section records episodic interactions (``user corrected my testing approach on 2026-02-14'').

\textbf{What compiles.} M encodes the user's \emph{type}---expertise-calibrated attention weighting, register adjustment (formal/casual/technical), empathetic framing priors. The model attends differently to the user's input: shorter latency to user intent, emphasis on aspects the user cares about, suppression of information the user already knows.

\textbf{What needs the slot.} Specific preference details (``prefers pytest over unittest''), names, dates, and particular past interactions.

\textbf{Why not a prompt?} Current approaches (system prompts, RAG-retrieved user profiles) place preferences in context where they (1) dilute over multi-turn conversation as the preference tokens become a shrinking fraction of total context, (2) incur per-turn token cost as the profile is re-injected each turn, and (3) are inspectable---the model can reference ``according to your profile'' rather than naturally adapting its behavior. Compiled M provides preference-shaped attention at zero token cost per turn. The Experiment 0 synthetic task (\S\ref{sec:synthetic_task}) validates this: 100\% Session 2 preference recall through M-state alone.

\textbf{Practical deployment.} An enmeshed assistant compiles the user's Mind Tree once per session (or on significant update). Subsequent turns benefit from preference-shaped attention without re-prompting. The M-state is swappable per user, enabling multi-tenant deployment from a single frozen host.

\subsection{Research Conviction and Intellectual Diversity}

\textbf{Mind Tree configuration.} The \texttt{beliefs} section encodes epistemic positions with categorical conviction labels. The \texttt{values} section encodes axiological commitments (``empirical evidence over theoretical elegance''). The \texttt{models} section encodes awareness of opposing positions (``mainstream consensus holds X; I disagree because Y'').

\textbf{What compiles.} M encodes \emph{dispositional conviction}---the tendency to defend positions under pressure, weight evidence in a particular direction, and resist consensus drift. This is the sycophancy resistance mechanism validated in \S\ref{sec:experiments}: F achieves 89\% resistance across 5 escalating pressure levels on the instruct model.

\textbf{What needs the slot.} Specific evidence, citations, statistical claims, and counter-arguments to predicted objections.

\textbf{Why not a prompt?} Research diversity in multi-agent systems requires agents that maintain genuinely distinct perspectives under debate pressure. Prompt-level beliefs are the easiest target for adversarial argumentation: the opposing agent can read the beliefs (or infer them from outputs) and construct targeted counter-arguments. When the beliefs are inspectable content, the model's chain-of-thought can reason about them and decide to abandon them. Compiled M is non-inspectable: the model cannot reason about its own Q/V modulation, so adversarial pressure targets the model's outputs rather than its dispositional anchoring. This produces agents that \emph{update on evidence} (the curated slot can be revised) while maintaining \emph{stable epistemic character} (the compiled disposition persists).

\textbf{Multi-agent implications.} In a research ensemble, each agent compiles a different Mind Tree encoding a distinct intellectual tradition, methodological commitment, or theoretical framework. The agents share the same frozen host but attend to different features of the evidence, extract different implications from shared data, and resist pressure to converge. Diversity is architectural rather than prompted.

\subsection{Agentic Strategic Thinking (ReAct / ALFWorld)}

\textbf{Mind Tree configuration.} The \texttt{strategies} section encodes procedural knowledge as scenario-based parent nodes with actionable child tactics. The \texttt{memories} section records episodic outcomes (``putting hot object on countertop failed---must cool first''). Strategy nodes carry \texttt{addresses} fields that map to task types (``find-and-heat'', ``examine-in-light'').

\textbf{What compiles.} M encodes \emph{procedural intuition}---approach tendencies (explore-vs-exploit), failure-pattern avoidance, and task-type recognition. The ALFWorld results (\S\ref{sec:experiments}) show the Mind Tree improving success from 3\% to 53\%, though the primary channel is structured context rather than compiled modulation for constrained action selection.

\textbf{What needs the slot.} Specific action sequences (``Step 1: go to fridge. Step 2: open fridge. Step 3: take apple''), tool signatures, and environment-specific details.

\textbf{Why not a prompt?} ReAct-style agents accumulate strategy logs, reflections, and episodic memories as text in context. Over multi-episode learning, this log grows linearly, eventually crowding out task-relevant context. Reflexion \citep{shinn2023reflexion} mitigates this with summarization but the summarized reflections still dilute. Compiled M encodes the \emph{distilled procedural tendency} at fixed cost: the agent's approach bias does not grow with experience. The strategy slot provides episode-specific tactics that may change each round, while M provides the stable procedural character that persists.

\textbf{Online learning loop.} After each episode, new observations are added to the Mind Tree (observation $\to$ tactic $\to$ strategy promotion via graph traversal). Re-compilation through $\phi$ updates M at the cost of one forward pass. The agent's procedural intuition evolves without retraining.

\subsection{Codebase and Repository Management}

\textbf{Mind Tree configuration.} The \texttt{strategies} section encodes architectural patterns (``this repo uses hexagonal architecture with ports-and-adapters''), coupling relationships (``auth middleware and session storage are tightly coupled''), and coding conventions. The \texttt{memories} section records past decisions and their rationale (``chose PostgreSQL over MongoDB for transactional consistency'').

\textbf{What compiles.} M encodes \emph{structural awareness}---which modules are likely coupled, where state lives, which patterns the codebase follows. The model navigates the codebase with shaped attention: when reading a file, M biases attention toward the patterns and structures that matter in this specific repository.

\textbf{What needs the slot.} Specific file paths, API signatures, dependency versions, and recent changes.

\textbf{Why not a prompt?} Repository context in prompt (README, architecture docs, file trees) scales linearly with codebase size. For large repositories, this context dominates the prompt budget and dilutes with conversation length. Compiled M provides structural intuition at $O(L \cdot d \cdot r)$ per token---constant in repository size. The curated slot injects only the file paths and API details relevant to the current query, selected by M's activation scores on strategy nodes. This produces a coding assistant that ``knows'' the repository structure implicitly rather than consulting it explicitly each turn.

\textbf{Multi-repository deployment.} Different codebases compile to different M-states. Switching repositories requires one forward pass through $\phi$, not re-prompting with a new architecture document. M-states can be cached per repository and loaded on demand.

\section{Enmeshed Network Design Space}
\label{sec:full_design_space}

We validate one point in a six-axis design space. Key unexplored directions:

\textbf{Gated blending (Level 2).} $h' = h + g(h) \odot f(h, M)$. A learned gate $g$ could solve the multi-turn attractor by learning when to suppress M---e.g., reducing modulation during conversational turns while maintaining it during identity-establishing moments.

\textbf{Cross-attention (Level 3).} $h' = h + \text{Attn}(h, M, M)$. Treats M as a specialized expert with attention as the implicit router. Could transfer richer content through the bottleneck but at higher computational cost ($O(L \cdot d^2)$ vs.\ $O(L \cdot d \cdot r)$).

\textbf{Adaptive rank allocation.} More capacity at layers where M's empirical perturbation magnitude is largest. Layer activity analysis shows bimodal distribution---some layers are consistently more active, suggesting non-uniform rank could optimize the capacity-cost tradeoff.

\textbf{Temporal profiles.} Decay, oscillation, or phase-gated M strength over generation steps. Constant M (this work) produces the attractor problem; a learned temporal profile could provide strong identity establishment followed by natural conversational engagement.

\textbf{Zone-aware modulation.} Different M scaling for belief zones vs.\ evidence zones vs.\ query zones within the context. Designed but not validated---preliminary experiments showed no significant effect over uniform scaling.

\textbf{Property graph backend.} The Mind Tree's hierarchical structure maps naturally to a property graph database with vector-indexed nodes, enabling multi-hop traversal, temporal belief tracking, and similarity-based discovery at scale beyond XML serialization.

\section{Loose Ends and Dead Ends}
\label{sec:loose_ends}

This section documents the research narrative---what we tried, what failed, and why---in the spirit of honest reporting that we hope saves other researchers time.

\textbf{Q-delta collapse (v10).} Contrastive loss on flattened M parameters pushed them apart (cosine sim 0.18--0.35) but the functional deltas $x \cdot M_A \cdot M_B^\top$ remained identical (cosine sim 0.9999). The hidden state $x$ dominates the product, washing out parameter-space differences. Functional contrastive loss on $\Delta Q$ directly was the fix.

\textbf{Epiphenomenal memory (v13.8).} The most surprising failure: 100\% retrieval accuracy, 100\% NTP wins, 0.508 nats improvement, measurable Q-perturbation (22\% magnitude, $28^\circ$ rotation)---and zero visible effect on generation. The model routed around M through the residual stream. This motivated causal necessity training (v14).

\textbf{The permission hypothesis.} M alone produces near-zero experiential output (3\% on the 4-condition benchmark). A system prompt alone produces generic hallucinations (43\%). Together: 93\%. M provides content; the system prompt provides ``permission'' to express it. This synergy was not designed---it emerged from the 4-condition factorial.

\textbf{Credence sensitivity.} Seven distinct training approaches failed to induce sensitivity to numeric credences (0.82 vs 0.45). The tokenizer representation bottleneck is fundamental: digit sequences are too similar in embedding space. Categorical labels (strong/moderate/agnostic) work because they use semantically distinct tokens.

\textbf{Bidirectional belief attention.} Allowing bidirectional attention within belief XML zones (while maintaining causal attention elsewhere) showed no significant effect. The model already attends well to beliefs with standard causal masking.

\textbf{Zone-aware M scaling.} Applying different M magnitudes to belief zones vs.\ evidence zones vs.\ query zones showed no improvement over uniform M application. The rank-16 bottleneck appears to be the binding constraint, not zone-specific magnitude.

\textbf{Compiled M for action selection.} Q/V modulation during ALFWorld generation hurts constrained action selection ($-13$ percentage points). M shifts token probabilities away from the exact action strings required. For constrained tasks, M's value is curation (selecting which strategies to put in context), not modulation.

\textbf{Environment-dependent PPL wins (v7.1).} The v7.1 checkpoint achieved 12/12 PPL wins only in the original training environment (Python 3.11 + older transformers). In a newer environment (Python 3.13 + transformers 4.56), base model forward pass produced different hidden states, breaking phi-base alignment. Lesson: pin your environment or expect to retrain phi.

\section{Reproducibility}
\label{sec:reproducibility}

Base model: Qwen3.5-4B-Base (HuggingFace). Training compute: single 24GB GPU, $\sim$10 hours total across all phases. Checkpoints and code will be released upon acceptance.

\textbf{Key hyperparameters:} rank $r = 16$, stride 3 (11 of 32 layers patched), $d = 2560$, max norm 1.5, belief window 512 tokens, v\_scale 1.0. Phase A: lr $10^{-4}$, margins 0.1--0.3. Phase B: lr $3 \times 10^{-4}$, margin 1.0. Phase D: lr $3 \times 10^{-5}$, 44 gold responses.
